\section{Related Work}
\label{sec:related}

\subsection{Geometric Shape Morphing}
Geometric approaches construct morphing trajectories from relationships
between the endpoint surfaces. Learned matching and interpolation methods
combine geometric features, optimal transport, and deformation priors to
estimate or refine correspondences without manual
annotations~\cite{eisenberger2020deepshells,eisenberger2021neuromorph,attaiki2023ncp}.
Neural implicit formulations further accommodate noisy inputs and topology
changes without supervision on intermediate shapes~\cite{sang2025fourdeform}.
These advances relax restrictive geometric assumptions, but primarily
address shape evolution rather than the joint synthesis of geometry and
appearance. Bridging geometric and generative approaches,
SRIF~\cite{sun2024srif} lifts multiview image morphs into intermediate point
clouds that guide continuous shape registration. Generated observations
can thus inform geometric trajectories, although jointly synthesizing
shape and texture also requires modeling how appearance evolves.

\subsection{Generative Image Morphing}
Diffusion priors allow morphing to synthesize intermediate content rather
than only deform and blend pixels. Existing approaches interpolate noisy
latents and conditioning signals, often together with endpoint-specific
adaptation or perceptually informed
sampling~\cite{wang2023interpolating,yang2024impus,zhang2024diffmorpher}.
Complementary strategies control feature mixing inside the generator
through attention interpolation, attention fusion, and
correspondence-guided
transport~\cite{he2024aid,cao2025freemorph,shen2024dreammover}.
Together, these works highlight the importance of coordinating global
semantic change with local structure and appearance, rather than relying
on latent interpolation alone. However, a smooth image transition does
not itself define a consistent 3D trajectory: independently lifting
intermediate images into 3D can introduce geometric discontinuities and
inconsistent unobserved surfaces. Transferring these mechanisms
to 3D therefore requires coupling their semantic flexibility with a
representation of the evolving asset.

\subsection{Generative Textured 3D Morphing}
Radiance fields and Gaussian primitives support detailed reconstruction
and rendering~\cite{mildenhall2020nerf,kerbl2023gaussians}, while learned
shape latents and structured Gaussian representations provide shared
spaces for generative
modeling~\cite{zhang2023shape2vecset,zhang2024gaussiancube}.
For textured morphing, volumetric transport and source-adapted diffusion
offer alternatives to explicit surface
deformation~\cite{tsai2022morphflow,Yang_2025_ICCV}.
The former optimizes transitions between multiview reconstructions,
whereas the latter interpolates noise, adapted model parameters, and
conditioning features within a pretrained 3D generator. Both construct
transitions using optimization specific to the input assets.
TRELLIS~\cite{xiang2025trellis} provides a versatile generative space
through structured latents (SLat), separating sparse spatial support from
local geometry and appearance features.

Building on this representation, recent training-free methods coordinate
geometry and texture through optimal transport, progressive alignment,
and attention-based
fusion~\cite{yin2025wukong,liu2026interp3d,sun2026morphany3d}.
WUKONG combines condition interpolation with sequential initialization
and semantic consistency, while Interp3D progressively aligns conditioning
semantics, SLat-guided structure, and appearance. MorphAny3D instead fuses
endpoint information through cross-attention and propagates preceding-frame
features through self-attention, with orientation correction to stabilize
the sequence. It also supports existing assets through 3D inversion.
Although these procedures adapt their outputs to each input pair, they do
not learn a shared morphing function across a collection of transitions.

\methodname{} takes a complementary approach by treating training-free
morphs as supervision rather than as the final inference procedure.
Using Interp3D-generated targets, we learn source fusion across many
object pairs and interpolation coefficients, conditioning both generation
stages directly on the source SLat representations. A separate frozen TRELLIS prior
regularizes the student's own outputs, complementing the example
trajectories supplied by the teacher. This yields a reusable
3D-conditioned morphing model that requires no source-specific parameter
optimization at inference.
